\documentclass{article}
\usepackage{amsmath, amsfonts, amsthm, amssymb}
\usepackage{enumitem}
\usepackage{graphicx} % Required for inserting images
\usepackage{geometry}
\usepackage[english]{babel}
\usepackage[affil-it]{authblk}
\usepackage[T1]{fontenc}
\usepackage[utf8]{inputenc}
\usepackage{lmodern}
\usepackage{tikz}
\usetikzlibrary{arrows.meta,positioning}
\usepackage{booktabs}
\usepackage{tabularx}
\usepackage{caption}
\usepackage{import}
\usepackage{bbm}
\title{Almost Disjointness Principles and $Q$-Space Cardinals}
\date{}
\theoremstyle{definition}
\newtheorem{definition}{Definition}
\newtheorem{prob}[definition]{Problem}
\newtheorem{thm}[definition]{Theorem}
\newtheorem{claim}{Claim}[definition]
\newtheorem{lemma}[definition]{Lemma}
\newtheorem{proposition}[definition]{Proposition}
\newtheorem{remark}[definition]{Remark}

\DeclareMathOperator{\dom}{dom}
\DeclareMathOperator{\Col}{Col}
\DeclareMathOperator{\htt}{ht}
\author{Vinicius de Oliveira Rodrigues\thanks{Electronic address: \texttt{vinior@ime.usp.br}}}
\affil{University of São Paulo}

\begin{document}\maketitle

\begin{abstract}
    Banakh and Bazylevych introduced separation-axiom variants $\mathfrak q_i$, for $i=1,2,2\frac{1}{2}$, of the cardinal $\mathfrak q$, together with a cardinal $\mathfrak{adp}$ lying between $\mathfrak{dp}$ and $\mathfrak{ap}$.
    They asked whether $\mathfrak{adp}$ coincides with either of these two cardinals.
    We prove in ZFC that $\mathfrak{adp}=\mathfrak{dp}$.
    We define a dual variant $\mathfrak{adp}_2$ and show that $\mathfrak{adp}_2=\mathfrak{ap}$.

    We further study the relation between $\mathfrak{ap}$ and the weakened $Q$-space cardinals.
    We introduce a tree analogue $\mathfrak{at}$ of $\mathfrak{ap}$ and prove $\mathfrak{at}=\mathfrak q_{2}$, concluding that $\mathfrak{ap}\leq\mathfrak q_{2}$, which answers a question of Banakh and Bazylevych.

    Assuming the Generalized Continuum Hypothesis, we construct ccc forcing extensions with $\mathfrak{ap}=\omega_1<\mathfrak{at}=\mathfrak c$, so $\mathfrak{ap}<\mathfrak q_2$ is consistent with ZFC.

    MSC: Primary 03E17, 03E05. Secondary 03E35, 54A35, 54D10.
\end{abstract}

\subimport{sections/}{introduction.tex}
\subimport{sections/}{adp.tex}
\subimport{sections/}{basics.tex}
\subimport{sections/}{diagonal.tex}
\subimport{sections/}{forcing.tex}
\subimport{sections/}{conclusion.tex}
\bibliographystyle{plain}
\bibliography{includes/biblio.bib}
\end{document}
